\chapter{Conclusiones}
\label{chap:conclusiones}

Esta tesis partió de una pregunta que parece sencilla: ¿qué información aporta un modelo probabilístico cuando una melodía se representa como una sucesión de clases de altura? El recorrido desde la cadena de Markov hasta el transformador mostró que la respuesta depende de aquello que se quiera conocer. Predecir la siguiente nota, proponer una trayectoria de estados y reducir el tiempo de cálculo son logros distintos. Ninguno de ellos basta, por sí solo, para ordenar los modelos en todos los usos posibles.

\section{Respuesta a la pregunta de investigación}

En la tarea que pudo evaluarse de forma común, la respuesta es clara. Con la fracción completa de entrenamiento, el transformador obtuvo una \emph{perplexity} de $5.823$, frente a $7.356$ del HMM finito, $7.968$ del HDP-HMM truncado y $8.022$ del VOMM. Como una \emph{perplexity} menor corresponde a una menor sorpresa ante las notas de prueba, el transformador fue el mejor predictor entre las configuraciones estudiadas. No se trata sólo de una diferencia en el promedio global. En la comparación por obra de la Tabla~\ref{tab:chapter4-pairwise-results}, la NLL del HMM finito superó a la del transformador en $0.2505$ nats por evento, con un intervalo del $95\%$ de $[0.2367,0.2641]$; los seis intervalos pareados excluyeron el cero. El resultado predictivo se sostiene, por tanto, tanto en el resumen global como al dar el mismo peso a cada obra.

Los dos resúmenes también advierten contra una lectura apresurada. La \emph{perplexity} global del HDP-HMM es ligeramente menor que la del VOMM, pero la comparación con igual peso por obra favorece al VOMM por $0.0336$ nats. La diferencia no es una contradicción: el primer cálculo da más peso a las piezas largas y el segundo considera cada obra como una unidad. Para la pregunta inferencial adoptada en esta tesis, el contraste pareado es el pertinente; el promedio global conserva su utilidad como descripción del conjunto de eventos.

La pregunta estructural recibe una respuesta más limitada. En los experimentos del Capítulo~\ref{chap:ihmm}, el HDP-HMM alcanzó una log-verosimilitud de ajuste mayor y utilizó menos estados que el HMM finito, tanto en el subconjunto clásico como en las $167$ obras comparables del corpus extendido. Esto muestra una representación latente más compacta dentro de esas corridas, pero no demuestra que sus estados correspondan a frases, acordes o secciones. La calidad estructural sigue sin evaluación externa. En consecuencia, el resultado principal de la tesis favorece al transformador para predecir clases de altura, pero no permite declarar un modelo superior para explicar la organización musical.

\section{El beneficio frente al costo}

Pese a su mayor complejidad, el transformador obtuvo un beneficio observable cuando el objetivo fue la predicción. Su ventaja apareció desde la fracción $0.25$ y aumentó al incorporar más obras, mientras que el HMM finito, el HDP-HMM y el VOMM mejoraron menos dentro del intervalo estudiado. Si se dispone de un corpus del tamaño utilizado, del mismo tipo de representación y de una GPU, esa ganancia justifica considerar el transformador aun cuando su arquitectura tenga muchos más parámetros. La conclusión no se extiende automáticamente a corpus pequeños: en la fracción $0.10$, los valores $7.640$ y $7.655$ del transformador y del HMM finito fueron demasiado cercanos para establecer una ventaja clara.

El tiempo observado tampoco permite identificar complejidad arquitectónica con lentitud. En la fracción completa, el transformador se ajustó en $169$ segundos, mientras que el HDP-HMM necesitó $578$ y el HMM finito $628$. Sin embargo, el primero se ejecutó en GPU y los otros en CPU. Estos tiempos describen la máquina y las implementaciones utilizadas; no prueban que el transformador sea intrínsecamente más rápido. El VOMM necesitó sólo $41$ segundos de ajuste, aunque su \emph{perplexity} fue mayor y su evaluación tomó más tiempo que la de los otros modelos. Como muestra la Tabla~\ref{tab:chapter4-cost-results}, no hubo un modelo que mejorara a todos los demás de manera simultánea en NLL, tiempo de ajuste y tiempo de evaluación.

Por ello, la pregunta de si vale la pena usar la tecnología más compleja sólo puede responderse en relación con un propósito. Para predecir la siguiente clase de altura en la escala estudiada, la mejora del transformador fue suficiente para justificar su uso bajo el equipo disponible. Si se necesita una referencia rápida y basada en conteos, el VOMM conserva una ventaja práctica en el ajuste. Si interesa examinar una trayectoria latente discreta, el HMM finito y el HDP-HMM proporcionan un objeto que el transformador no ofrece de forma equivalente: el primero predijo mejor en la corrida final y el segundo produjo representaciones más compactas en los experimentos descriptivos. Esta última propiedad sólo constituye una ventaja si la compacidad es relevante para el análisis; todavía no equivale a interpretabilidad musical.

La reproducibilidad forma parte del mismo balance. Las particiones por obra, las semillas, la auditoría de cobertura, la caché y los puntos de reanudación permitieron terminar y revisar una corrida larga. A la vez, la falta de los CSV y JSON originales impide una auditoría tan directa como la que permitirían los artefactos primarios. Una tecnología no resulta conveniente sólo porque alcance la mejor métrica: también debe poder reconstruirse con los recursos y registros disponibles.

\section{Alcance y pasos siguientes}

La conclusión queda acotada por la representación utilizada. Reducir la partitura a clases de altura permitió comparar cuatro familias en la misma tarea, pero dejó fuera duración, registro, dinámica y buena parte de la organización métrica. El transformador asignó mejores probabilidades a la siguiente clase de altura; no se ha establecido cómo cambiaría la comparación al recuperar esas otras dimensiones.

El siguiente paso decisivo es incorporar anotaciones estructurales independientes en las mismas obras. Con ellas podrían evaluarse fronteras y particiones mediante $F_1$, información mutua normalizada e índice de Rand ajustado. Sólo entonces sería posible saber si la compacidad del HDP-HMM o las trayectorias del HMM finito guardan relación con una lectura musical reconocible. Hasta realizar esa comparación, los estados latentes y los patrones de atención deben entenderse como hipótesis de organización, no como análisis musicales validados.

También quedan comprobaciones experimentales concretas. Es necesario ampliar la rejilla del HMM finito más allá de $48$ estados, repetir la partición con otra semilla, comparar los tiempos sobre recursos equivalentes y recuperar los artefactos primarios de la corrida. Estas tareas permitirían distinguir entre el límite de una configuración y el de una familia de modelos, además de precisar el costo real de reproducir cada alternativa.

La aportación central no es, entonces, una clasificación universal de tecnologías. Es una forma de plantear la comparación sin confundir objetivos. Bajo la tarea común y la escala observada, el transformador ofrece el beneficio predictivo más claro. Para describir estados latentes, estudiar su compacidad o trabajar con un presupuesto distinto, los modelos probabilísticos más simples conservan funciones propias. Elegir entre ellos requiere decidir primero qué se desea explicar, qué evidencia permite comprobarlo y qué costo puede sostenerse de manera reproducible.
